\chapter{Methodology}\label{chap:methodology}

\begin{itemize}
    \item Mutual information, information theory:
    \begin{align*}
        I(X,Y)=H(X)+H(Y)-H(X,Y)
    \end{align*}
    where $H$ is the entropy (measure of complexity)
    \item Green matrix, Min-/Meet-matrix
    \item do multiple samples and see if critical $\epsilon$ changes (stability under resampling)
\end{itemize}

For two simplicial complexes $K_1$ and $K_2$, we define a scalar product as 
\begin{align*}
    \langle K_1,K_2\rangle_f&:=f(K_1)+f(K_2)-f(K_1\cup K_2),
\end{align*}
where $f$ is any nonnegative integer-valued (subadditive) function defined on simplicial complexes. For example $f(K)=\text{cat}(K)$. Due to the subadditivity of $f$, the scalar product is nonnegative. A stronger notion would be submodularity. Then it would hold $\left\langle K_1,K_2\right\rangle_f\ge f(K_1\cap K_2)\ge0$.

\todo[inline]{add discussion about $\langle K_1,K_2\rangle_f=0\iff$??}

For a finite sequence of filtration values $\epsilon_1,\dots,\epsilon_N$ we can define a matrix $\Sigma^{(N)}=(\Sigma^{(N)}_{ij})$ by
\begin{align*}
    \Sigma^{(N)}_{ij}:=\left\langle K_{\epsilon_i}(X),K_{\epsilon_j}(X)\right\rangle_f,\quad i,j=1,\dots,N.
\end{align*}
Since we consider a filtration $K_{\epsilon_1}\subseteq\dots\subseteq K_{\epsilon_{N}}$, the scalar product is just given by $\left\langle K_{\epsilon_i},K_{\epsilon_j}\right\rangle_f=f(K_{\epsilon_i}\cap K_{\epsilon_j})=f(K_{\epsilon_{i\wedge j}})$, where $i\wedge j:=\min\{i,j\}$. This means that the subadditivity of $f$ is not needed for the nonnegativity of the scalar product when considering a filtration. Setting $c_i=f(K_{\epsilon_i})$, the scale similarity matrix has the following form:
\begin{align}\label{eq:Sigma_form}
    \Sigma^{(N)}=\begin{pmatrix}
        c_1 & c_1 & \dots & c_1\\
        c_1 & c_2 & \dots & c_2\\
        \vdots & \vdots & \ddots & \vdots\\
        c_1 & c_2 & \dots & c_N
    \end{pmatrix}
\end{align}

The following lemma shows that the local extrema of the vector $c=(c_1,\dots,c_N)$ and the principal eigenvector $v^*$ of $\Sigma^{(N)}$ coincide. In particular, choosing a functional $f$ such that $i\mapsto f(K_{\epsilon_i})$ is monotone might not be a good idea, since the entries of the principal eigenvector will merely reflect the monotone behavior of $f$. This suggests that we need to choose a functional $f$ such that $i\mapsto f(K_{\epsilon_i})$ is non-monotone.

% The following lemma shows that using the $\min$ kernel might not be a good idea, as we do not gain any advantages from computing the principal eigenvector $v^*$ of $\Sigma^{(N)}$. Concretely, the entries of any eigenvector of $\Sigma^{(N)}$ have the same behavior as the map $i\mapsto f(K_{{\epsilon_i}})$.

\begin{Lem}
    Let $\Sigma\in\mathbb{R}^{N\times N}$ be defined by $\Sigma_{ij}=f(K_{\epsilon_{i\wedge j}})$ and denote $c_i:=f(K_{\epsilon_i})$ for $i=1,\dots,N$. Then for the principal eigenvector $v$ of $\Sigma$ it holds that
    \begin{align*}
        \operatorname{sgn}(v_{i+1}-v_i)=\operatorname{sgn}(c_{i+1}-c_i),\quad\forall i=1,\dots,N-1.
    \end{align*}
\end{Lem}
\begin{proof}
    Let $\lambda$ be the largest eigenvalue of $\Sigma$ with corresponding (principal) eigenvector $v$. Then $\Sigma v=\lambda v$. By Perron-Frobenius $\lambda>0$. We remember that the matrix $\Sigma$ has the form~\eqref{eq:Sigma_form} and compute
    \begin{align*}
        \lambda(v_{i+1}-v_i)&=\lambda v_{i+1}-\lambda v_i\\&
        =(\Sigma v)_{i+1}-(\Sigma v)_i\\&
        =\left(c_1v_1+\dots+c_{i}v_{i}+c_{i+1}\sum_{l=i+1}^{N}v_l\right)\\&-\left(c_1v_1+\dots+c_{i-1}v_{i-1}+c_i\sum_{l=i}^{N}v_l\right)\\&
        =c_{i+1}\sum_{l=i+1}^{N}v_l-c_i\sum_{l=i+1}^{N}v_l\\&
        =(c_{i+1}-c_i)\sum_{l=i+1}^{N}v_l.
    \end{align*}
    Since by Perron-Frobenius all entries of $v$ are positive, we have $\sum_{l=i+1}^{N}v_l>0$ in addition to $\lambda>0$. It follows that $\text{sgn}(v_{i+1}-v_i)=\text{sgn}(c_{i+1}-c_i)$ which is the claim.
\end{proof}

An interesting question arises concerning the global extrema of $c$ and $v^*$. Do they coincide as well? If not, how would counterexamples look like?
\todo[inline]{add discussion about global extrema}

\begin{Prop}[Properties of the Min-matrix]\label{prop:properties_of_min_matrix}
    Let $\Sigma=(\Sigma_{ij})\in\mathbb{R}^{N\times N}$ with $\Sigma_{ij}=\min\{c_i,c_j\}$ for some constants $c_1,\dots,c_N\in\mathbb{R}_{\ge0}$.
    \begin{enumerate}
        \item[1.] $\Sigma$ is positive semi-definite iff $0\le c_1\le\dots\le c_N$
        \item[2.] $\Sigma$ is positive definite iff $0<c_1<\dots<c_N$
        \item[3.] The determinant of $\Sigma$ is given by $\operatorname{det}(\Sigma)=c_1(c_2-c_1)\cdots(c_N-c_{N-1})$
        \item[4.] The inverse of $\Sigma$ exists iff $c_{i+1}-c_{i}\ne0$ for all $i=1,\dots,N-1$ and is given by $\Sigma^{-1}=(\Sigma^{-1}_{ij})$ with
        \begin{align*}
            \Sigma^{-1}_{ii}&=\begin{cases}
                \frac{1}{c_i-c_{i-1}}+\frac{1}{c_{i+1}-c_i} & ,i=1,\dots,N-1\\
                \frac{1}{c_{N}-c_{N-1}} & ,i=N
            \end{cases}\\
            \Sigma^{-1}_{i,i+1}&=\Sigma^{-1}_{i+1,i}=-\frac{1}{c_{i+1}-c_i},i=1,\dots,N-1
        \end{align*}
    \end{enumerate}
\end{Prop}
\begin{proof}
    \todo[inline]{Prove in appendix or cite}
\end{proof}

% Question: under which conditions on $f$ is $\Sigma$ positive definite or positive semi-definite

\begin{Rem}
    From Proposition~\ref{prop:properties_of_min_matrix} it follows that the scale similarity matrix $\Sigma$ is positive definite iff $i\mapsto f(K_{\epsilon_i})$ is monotonously strictly increasing and positive.
\end{Rem}

A matrix $\Sigma=(\Sigma_{ij})\in\mathbb{R}^{N\times N}$ is called nonnegative (positive) if $\Sigma_{ij}\ge0$ ($\Sigma_{ij}>0$) for all $i,j=1,\dots,N$. It is called primitive if it is nonnegative and there exists a natural number $m$ such that $\Sigma^{m}$ is positive. The spectrum $\sigma(\Sigma)$ is the (multi-)set of all eigenvalues of $\Sigma$: $\sigma(\Sigma)=\{\lambda\in\mathbb{C}:\exists v\in\mathbb{C}^N,v\ne0:\Sigma v=\lambda v\}$.

\todo[inline]{add definition of irreducible matrix}
\todo[inline]{discuss conditions on matrices}

% \todo[inline]{algebraic multiplicity = 1 $\implies$ geometric multiplicity = 1 $\implies$ eigenvector is unique up to scaling}

\begin{Thm}[Perron-Frobenius]
    Let $\Sigma=(\Sigma_{ij})\in\mathbb{R}^{N\times N}$ be a primitive matrix. Then there exist $\lambda^*>0$ and $v^*=(v^*_1,\dots,v^*_N)\in\mathbb{R}^{N}_{>0}$ such that $\lambda^*\in\sigma(\Sigma)$, $\lambda^*>|\lambda|$ for all $\lambda\in\sigma(\Sigma),\lambda\ne\lambda^*$. Furthermore, $v^*$ is the corresponding eigenvector that is unique up to scaling.
\end{Thm}

\begin{Rem}
    Analogously to classical PCA, we can study the spectral gap ratio: $\frac{\lambda_N-|\lambda_N|}{\lambda_N}$
\end{Rem}

\subsubsection*{critical scale choice}

We are now interested in $j_{\text{crit}}=\text{argmax}_{j=1,\dots,N}\{v^{\star}_j\}$. We pick $\epsilon_{\text{crit}}=\epsilon_{j_{\text{crit}}}$ and consider all vertices that are isolated in $\text{VR}_{\epsilon_{\text{crit}}}(X)$ as noise.

\todo[inline]{Discuss uniqueness of $j_{\text{crit}}$}

\subsubsection*{Connection to Markov Chains}

\todo[inline]{Doob $h$-transform}

Let $\Sigma_N$ be the scale similarity matrix with principal eigenvector $v^*=(v^*_1,\dots,v^*_N)$ and corresponding eigenvalue $\lambda^*$. Let $\{1,\dots,N\}$ be the state space and introduce the transition matrix
\begin{align*}
    P=\frac{1}{\lambda^*}D_{v^*}^{-1}\Sigma_ND_{v^*},
\end{align*}
where $D_{v^*}=\operatorname{diag}(v^*_1,\dots,v^*_N)$. The resulting Markov chain is reversible and ergodic. Its stationary distribution $\pi=(\pi_1,\dots,\pi_N)$ is given by
\begin{align*}
    \pi_j=\frac{v^*_j}{\sum_{k=1}^{N}(v^*_k)^2}
\end{align*}

\begin{Rem}
    The Donsker-Varadhan free-energy functional $\mathcal{F}$ is known to be
    \begin{align*}
        \mathcal{F}_{\Sigma}(Q)=\sum_{i,j=1}^{N}\pi_i^{Q}q_{ij}\log{\left(\frac{\Sigma_{ij}}{q_{ij}}\right)}
    \end{align*}
    where $Q=(q_{ij})$ is any positive, row-stochastic matrix with stationary distribution $\pi^{Q}$. It holds $\mathcal{F}(Q)\le\mathcal{F}(P)$ for any positive, row-stochastic matrix $Q$, where $P$ is the matrix introduced in Equation~[]. This fact serves as motivation for the replacement of our initial scale similarity matrix $\Sigma_N$ with the transition matrix $P$: we replace $\Sigma_N$ which is non-stochastic by the closest (in the above sense) stochastic matrix $P$.
\end{Rem}

\section{Monotone scale selection}

We consider the case $0<c_1<\dots<c_N$. Define the discrete Laplacian:
\begin{align*}
    \mathcal{L}=\frac{1}{c_1}\operatorname{diag}(1,0,\dots,0)-\Sigma^{-1}
\end{align*}
which acts as the generator of a birth-death chain on the state space $\{1,\dots,N\}$ with transition rates...

The normalized stationary distribution of the embedded birth-death chain is given by
\begin{align*}
    \pi^{\text{BD}}&=(\pi^{\text{BD}}_1,\dots,\pi^{\text{BD}}_N)\\
    \pi^{\text{BD}}_j&=\frac{r_j}{\overline{r}}\\
    r_j&=\begin{cases}
        \frac{1}{c_2-c_1}, & j=1\\
        \frac{1}{c_j-c_{j-1}}+\frac{1}{c_{j+1}-{c_j}}, & j=2,\dots,N-1\\
        \frac{1}{c_N-c_{N-1}}, & j=N
    \end{cases}\\
    \overline{r}&=\sum_{j=1}^{N}r_j
\end{align*}

The critical scales are naturally identified as
\begin{align*}
    j_{\text{crit}}\in\arg\min\{\pi^{\text{BD}}_j:j=1,\dots,N\}
\end{align*}

\section{Other combinatorial and topological features of simplicial complexes}

In this section we present and discuss function defined on simplicial complexes that are non-negative, integer valued and sub-additive. Let $\mathcal{F}(K)$ denote the set of all facets of $K$, i.e. all simplices $\sigma\in K$ with $\dim\sigma=\dim K$.

\textbf{$f$-vector.} The $f$-vector of a simplicial complex $K$ is $f(K)=(f_0,f_1,\dots,f_d)$, where $f_k$, $k=-1,0,\dots,d$ is the number of simplices in $K$ containing $k+1$ vertices. Using the $f$-vector, we can construct the $f$-polynomial:
\begin{align*}
    F_K(t)=\sum_{k=-1}^{d}f_kt^{k+1}=1+f_0t+f_1t^2+\dots+f_dt^{d+1}
\end{align*}

\textbf{Betti numbers.} Informally speaking, the $k$-th Betti number $\beta_k(K)$ of a simplicial complex $K$ represents the number of $k$-dimensional holes. In particular, $\beta_0(K)$ is the number of connected components of $K$. We could also consider the total Betti number (the sum of all Betti numbers): $\sum_{k}\beta_k(K)$.

\textbf{Euler characteristic.} The Euler characteristic is given by
\todo{motivation, proof}
\begin{align*}
    \chi(K)=\sum_{k=0}^{d}(-1)^{k}f_k(K)=\sum_{k=0}^{d}(-1)^{k}\beta_k(K)
\end{align*}

\textbf{Lusternik-Schnirelmann Category.} 

\textbf{Vertex cover number.} The vertex cover number is defined as
\begin{align*}
    \tau(K):=\min\{|C|:C\subseteq V(K),C\cap F\ne\emptyset\forall F\in\mathcal{F}(K)\}
\end{align*}

\textbf{Matching number.} The matching number is defined as
\begin{align*}
    \nu(K):=\max\{|M|:M\subseteq\mathcal{F}(K),F\cap F'=\emptyset\forall F,F'\in M, F\ne F'\}
\end{align*}

The matching number can be computed using the blossom algorithm.

\textbf{Chain count.} The chain count of length $r\in\mathbb{N}$ is defined as
\begin{align*}
    c_r(K):=|\{(\sigma_0\nsubseteq\sigma_1\nsubseteq\dots\nsubseteq\sigma_r):\sigma_i\in K\}|
\end{align*}

\textbf{Degree-type counts.} The degree-type counts of dimension $t\in\mathbb{N}_0$ are defined as
\begin{align*}
    \Delta_t(K):=\max_{\tau\in K}|\{\sigma\in K:\tau\subseteq\sigma,\dim\sigma=t\}|
\end{align*}

% \begin{itemize}
%     \item graph based descriptors of the 1-skeleton (graph-theoretic invariants): vertex cover number, number of connected components, arboricity, minimum spanning tree
%     \item general construction?
%     \item instead of $f$ vector count $k$ simplexes that are not part of $k+1$ simplexes
%     \item integer linear programming?
%     \item use persistence diagrams (distance between persistence diagrams), bar codes, birth and death intervals to construct matrix (persistence landscape norm, persistence entropy)
%     \item $n_k(\epsilon)=\#\{(b,d)\in\text{PD}_k:b\le\epsilon<d\}$ or $L_k(\epsilon)=\sum_{(b,d)\in\text{PD}_k,b\le\epsilon<d}(d-b)$
% \end{itemize}

\begin{Rem}
    Here we discuss how the functionals presented above simplify when the sampled points live on a line $\mathbb{R}^1$. Then there is only the $0$-th Betti number $\beta_0$ (the number of connected components). The $f$-vector only has one non-trivial component: the number of edges $f_1$.
\end{Rem}


We consider the following general form of the functional:
\begin{align*}
    f(K)=f_{\text{Top}}(K)\cdot f_{\text{Comb}}(K),
\end{align*}
where the topological part $f_{\text{Top}}$ is defined by
% \begin{align*}
%     f_{\text{Top}}(K)=\delta_{-1}+\sum_{j=0}^{d-1}\delta_j\beta_j^{\alpha_j}(K)+\delta_d\operatorname{cup}_R^{\alpha_d}(K)+\delta_{d+1}\operatorname{PC}^{\alpha_{d+1}}(K)
% \end{align*}
\begin{align*}
    f_{\text{Top}}(K)=\delta_{-1}+\sum_{j=0}^{d-1}\delta_j\beta_j^{\alpha_j}(K)
\end{align*}
with coefficients $\delta_j\in\mathbb{R}$, $j=-1,0,\dots,d-1$ and $\alpha_j\in\mathbb{R}$, $j=0,\dots,d-1$. The combinatorial part $f_{\text{Comb}}$ is defined by
% \begin{align*}
%     f_{\text{Comb}}(K)&=\gamma_{-1}+\sum_{j=0}^{d}\gamma_jf_j^{\xi_j}(K)+\gamma_{d+1}\operatorname{CR}_{k,m}^{\xi_{d+1}}(K)+\gamma_{d+2}c_r^{\xi_{d+2}}(K)\\&
%     +\gamma_{d+3}\Delta_t^{\xi_{d+3}}(K)+\gamma_{d+4}\tau^{\xi_{d+4}}(K)+\gamma_{d+5}\nu^{\xi_{d+5}}(K)
% \end{align*}
\begin{align*}
    f_{\text{Comb}}(K)&=\gamma_{-1}+\sum_{j=0}^{d}\gamma_jf_j^{\xi_j}(K)
\end{align*}
with coefficients $\gamma_j\in\mathbb{R}$, $j=-1,0,\dots,d$ and $\xi_{j}\in\mathbb{R}$, $j=0,\dots,d$.

This function serves as a building kit that lets us construct a polynomial-type combination of all the features that we introduced earlier. For example, setting $\delta_0=1$, $\delta_j=0$ for $j\ne 0$, $\gamma_1=1$, $\gamma_j=0$ for $j\ne1$ and $\alpha_j=\xi_j=1$ recovers the product $\beta_0f_1$ which is the number of edges $f_1$ multiplied by the number of connected components $\beta_0$ (see \cite{leonov-report}). If we consider a filtration $K_{\epsilon}$, then if $\epsilon$ increases, $f_1$ increases (more edges appear), but $\beta_0$ decreases (connected components merge), so this is a non-monotone kind of functional.

\todo[inline]{Talk about how to recover the Euler characteristic using this functional}


\begin{Def}
    A triple $(f,\{K_i\},\oplus)$ is said to be spectrally optimizable if it generates a finite irreducible discrete time Markov chain on the state space $\{1,\dots,N\}$ with transition probability matrix $Q=(q_{ij})$ and stationary distribution $\pi^Q=(\pi_1^Q,\dots,\pi_N^Q)$.
\end{Def}